\section{Applications}
\label{sec:applications}

\subsection{Calculating differential properties}

The proposed representation of the difference distribution table allows for the efficient computation of cryptographic metrics without the need to exhaustively generate the standard tabular format. In the following, we demonstrate how each of the differential properties introduced previously can be calculated using this data structure.

To calculate the differential uniformity, it is sufficient to inspect the values stored at the terminal nodes of the data structure. Let $T$ be the set of all terminal nodes. Since the differential uniformity is the maximum frequency of any non-trivial differential, it corresponds directly to the maximum integer value among these nodes. Finding this maximum requires a simple traversal, which takes $\mathcal{O}(|T|)$ time. Since the number of terminal nodes is at most $\Delta_{F}$, the cost is negligible compared to the size of the original DDT.

To compute the differential spectrum, we traverse the automaton starting from the root using a standard search algorithm, such as depth-first search. During the traversal, we track the multiplicity of the current path by multiplying the number of valid transition symbols along each traversed edge. When a terminal node storing a value $i$ is reached, the accumulated path weight is added to the corresponding counter $\omega_i$. This approach requires $\mathcal{O}(D)$ memory, where $D$ is the depth of the automaton and is bounded by the input size $n$. The time complexity is $\mathcal{O}(P)$, where $P$ is the total number of distinct valid paths in the reduced structure.

Impossible differentials, which correspond to zero-entries in the difference distribution table, can also be efficiently extracted from the data structure. To determine if a differential $(a, b)$ is impossible, it suffices to evaluate its encoding, where symbols are formed by the respective bits of $a$ and $b$. That is, if $a = 101$ and $b = 011$, the encoding of the differential $(a, b)$ is the symbol sequence $213$, since the bit pairs from most to least significant are $(1,0)$, $(0,1)$, and $(1,1)$. Since the automaton only accepts non-zero entries, the differential is impossible if and only if its encoding is not in $L(DDT)$. However, for verifying a single specific differential, this traversal is not computationally advantageous compared to a standard DDT. While a direct table lookup operates in $\mathcal{O}(1)$ time, evaluating through the automaton requires $\mathcal{O}(D)$ operations. While this overhead makes the automaton less efficient than a DDT for single-point evaluations, the $\mathcal{O}(D)$ cost remains bounded by the input size $n$ and may not represent a concern. Next, we present methods for performing more complex queries on the automaton, enabling the simultaneous extraction of multiple impossible differentials and other properties.

\subsection{Complex Queries via Automata Operations}
\label{sec:queries}

The transformation of the DDT into a finite automaton establishes a formal
language containing the data in the table. By treating the set of valid
differentials as a language, complex cryptanalytic queries can be
formulated as regular expressions and resolved as operations on this
language. This transforms the ad-hoc process of searching for specific
differential properties into an automated, language-theoretic problem.
Formally, the query process follows these steps:
\begin{enumerate}
    \item The query is formulated as a regular expression $R$ over the
          alphabet $\Sigma$.
    \item The regular expression $R$ is compiled into a finite automaton
          $Q$.
    \item The intersection of $D$ and $Q$ is computed using methods such
          as topological dynamic programming with bit-parallelism or lazy
          depth-first search.
    \item The resulting intersection language is then enumerated, yielding
          the set of DDT entries that satisfy the properties specified by
          the regular expression $R$.
\end{enumerate}

This formal language approach provides a highly flexible mechanism. Rather
than developing custom algorithms for specific search criteria, an analyst
simply constructs the appropriate regular expression. Moreover, because
the class of regular languages is closed under union, intersection, and
complement, complex queries can be assembled compositionally from simpler
ones without leaving the framework. The same query automata can be combined
with the complement of $D$ to reason about impossible differentials, with a
value-restricted sub-automaton of $D$ to filter by probability, or with each
other to express conjunctive constraints.

Throughout the rest of this section, we write the input difference as
$a = a_{n-1}\cdots a_0$ and the output difference as $b = b_{n-1}\cdots
b_0$. Following the encoding of Section~3, each symbol of the word
representing $(a,b)$ is built from a pair of bits $(a_i, b_i)$, so the
four symbols partition naturally into two pairs: symbols
$\{\texttt{0},\texttt{2}\}$ correspond to $a_i = 0$ (input bit
inactive) and symbols $\{\texttt{1},\texttt{3}\}$ to $a_i = 1$ (input
bit active), while symbols $\{\texttt{0},\texttt{1}\}$ correspond to
$b_i = 0$ (output bit inactive) and symbols $\{\texttt{2},\texttt{3}\}$
to $b_i = 1$ (output bit active). For brevity, we display regular
expressions in typewriter font and adopt the usual convention that the
dot \texttt{.} matches any symbol of $\Sigma$, that is, \texttt{.}
$\equiv$ \texttt{(0|1|2|3)}. To illustrate the practical utility of
this method, the following examples demonstrate how common queries can
be formulated for a DDT of size $2^n$.

\begin{itemize}
    \item Fixed input differences: To identify all possible output
          differences for a specific $n$-bit input difference
          $a = \texttt{0101}$, the regular expression is formulated as
          \[
              R = \texttt{(0|2)(1|3)(0|2)(1|3)}.
          \]
          This query replaces the conventional $O(2^n)$ row scan required
          in a standard DDT. Furthermore, it can be reused for other
          values of $a$ by simply changing the transition symbols in the
          query automaton.

    \item Partial and truncated differentials: To evaluate a truncated
          differential where only the two most significant input bits are
          fixed to $\texttt{10}$, the regular expression is formulated as
          \[
              R = \texttt{(1|3)(0|2)..}
          \]
          More generally, any pattern in $\{0,1,*\}^{n}$ on input bits and
          $\{0,1,*\}^{n}$ on output bits maps to a length-$n$ regular
          expression with at most four alternatives per position, so that
          any truncated query of arbitrary mask shape compiles into an
          automaton of $O(n)$ states. The special case in which each
          position is constrained only to be active or inactive (rather
          than fixed to a specific bit value) recovers the activity
          patterns that underlie the MILP- and SAT-based active-S-box
          counts of Mouha et al.~\cite{MouhaEtAl2012} and the differential and linear
          bounds of ciphers such as SKINNY~\cite{BeierleEtAl2016}.

    \item Impossible differentials: Impossible differentials are precisely
          the words rejected by $D$. Rather than testing individual pairs
          $(a,b)$, the framework exposes the entire impossible language
          as the complement
          \[
              L(I) = \overline{L(D)} \cap \Sigma^{n},
          \]
          which is recognized by the complement automaton $\overline{D}$.
          For instance, to enumerate all impossible differentials of a
          $4$-bit S-box whose input has only the least significant bit
          active ($a = \texttt{0001}$, $b$ unconstrained), the regular
          expression is
          \[
              R = \texttt{(0|2)(0|2)(0|2)(1|3)},
          \]
          and the impossible cells of that row are read from
          $\overline{L(D)} \cap L(R)$. This is the elementary
          building block of impossible-differential cryptanalysis as
          first applied to Skipjack~\cite{BihamEtAl1999} and DEAL~\cite{Knudsen1998}, where one
          routinely needs the full set of zero-entries that respect a
          given activity pattern.

    \item Extremal-probability differentials: Because the reduced
          automaton stores the differential multiplicity directly on
          each accepting state, filtering by probability does not
          require any path enumeration: it is a one-pass scan of the
          state set. Let $D_{\geq \tau}$ denote the sub-automaton
          obtained from $D$ by retaining only the accepting states
          whose stored value is at least $\tau$ and pruning unreachable
          states. Then $L(D_{\geq \tau})$ is the set of differentials
          with probability at least $\tau \cdot 2^{-n}$, and the
          differentials achieving the differential uniformity
          $\Delta(F)$ are obtained by setting $\tau = \Delta(F)$. For
          the PRESENT S-box~\cite{BogdanovEtAl2007}, where $n = 4$ and $\Delta = 4$, the
          language $L(D_{\geq 4})$ enumerates the handful of
          high-probability transitions that anchor the best known
          characteristics of the cipher; choosing $\tau = 2$ instead
          returns every non-trivial valid transition that is available
          to a trail. This is exactly the seed information consumed by
          branch-and-bound differential trail searches such as
          Matsui's~\cite{Matsui1995} and its modern MILP and SAT descendants. The
          construction of $D_{\geq \tau}$ visits each accepting state
          at most once, so it runs in $O(|S|)$ time independently of
          the size of the resulting language.

    \item Hamming-weight bounded differentials: Many design strategies and
          search heuristics rely on input and output differences of low
          Hamming weight, since these often dominate the best multi-round
          trails. The set of words encoding differentials with
          $\mathrm{wt}(a) + \mathrm{wt}(b) \leq w$ is recognized by a
          counter automaton with $O(n\,w)$ states: at each position, the
          automaton increments its weight counter by $0$, $1$, or $2$
          depending on whether the symbol is $\texttt{0}$, in
          $\{\texttt{1},\texttt{2}\}$, or $\texttt{3}$.
          For example, restricting to single-bit input and output
          ($\mathrm{wt}(a) = \mathrm{wt}(b) = 1$) is recognized by the
          regular expression
          \[
              R = \texttt{0}^{*}\texttt{2}\,\texttt{0}^{*}\texttt{1}\,\texttt{0}^{*}\;|\;
                  \texttt{0}^{*}\texttt{1}\,\texttt{0}^{*}\texttt{2}\,\texttt{0}^{*}\;|\;
                  \texttt{0}^{*}\texttt{3}\,\texttt{0}^{*},
          \]
          and its intersection with $D$ enumerates the differentials
          targeted by single-bit and bit-flip distinguishers.
          Intersecting more general counter automata with $D$ supports
          the wide-trail strategy of Daemen and Rijmen~\cite{DaemenRijmen2002}, in which
          the designer must verify that no low-weight differential of a
          chosen S-box falls below the targeted probability threshold.

    \item Linearly constrained differentials: Linear relations on the
          bits of $(a,b)$ arise naturally when an S-box is embedded in
          a larger primitive. In a multi-round setting, the linear
          layer that follows a round of S-boxes constrains how the
          output differences of one round map to the input differences
          of the next, which translates, after fixing the surrounding
          differences, into a linear constraint on each S-box's local
          $(a,b)$. The simplest such constraint is the diagonal
          relation $a = b$, which is recognized by the regular
          expression
          \[
              R = \texttt{(0|3)}^{n},
          \]
          since symbols $\texttt{0}$ and $\texttt{3}$ are exactly those
          with $a_i = b_i$.
          More generally, any $\mathbb{F}_{2}$-linear relation on the
          bits of $(a,b)$ is recognized by a deterministic automaton
          whose states track the running parities of the relation, with
          at most $2^{r}$ states for $r$ independent linear constraints.
          Intersecting such an automaton with $D$ extracts exactly the
          DDT entries that are compatible with the linear context, a
          search that traditionally requires either bespoke filters or
          auxiliary MILP or SAT formulations.

    \item Composite queries: Because regular languages are closed under
          boolean operations, the previous queries can be combined freely.
          A concrete example is the search for high-probability
          low-weight differentials, formulated as the intersection
          \[
              L(D_{\geq \Delta(F)/2}) \cap L(\mathrm{wt} \leq 2),
          \]
          which is precisely the kind of seed used by trail-search
          heuristics for ciphers such as PRESENT~\cite{BogdanovEtAl2007} and SKINNY~\cite{BeierleEtAl2016},
          where only differentials that are both probable enough to
          contribute and sparse enough to survive the diffusion layer
          need to be retained. Other useful combinations include the
          intersection of $\overline{D}$ with a truncated query, which
          enumerates activity patterns guaranteed to vanish across one
          S-box layer, and the intersection of a linear-constraint
          automaton with a truncated query, which filters
          the differentials that survive a linear layer while
          respecting a required activity pattern. Each compound query
          is built by a single automata product and inherits the same
          enumeration and counting primitives as the elementary ones.

    \item Comparing two DDTs: The framework also supports queries that
          span two S-boxes. Given automata $D_F$ and $D_G$ for the DDTs
          of $F$ and $G$, the language $L(D_F) \setminus L(D_G)$
          collects the differentials that are possible for $F$ but
          impossible for $G$, and $L(D_F) \cap L(D_G)$ collects those
          shared by both. Such operations are useful for comparing
          candidate S-boxes during the design of a cipher: members of
          the same affine-equivalence class share the same DDT up to a
          relabelling of inputs and outputs, so $L(D_F) \setminus
          L(D_G)$ provides a direct certificate of inequivalence,
          while comparing two candidates from different equivalence
          classes (a situation that arises, for example, in the S-box
          selection of lightweight ciphers such as PRINCE~\cite{BorghoffEtAl2012} and
          Midori~\cite{BanikEtAl2015}) isolates the differentials that distinguish
          their behaviour under a fixed differential property of
          interest.
\end{itemize}

For all of the queries above, the running time is dominated by the size
of the intersection language, which is bounded by $|D| \cdot |Q|$ in the
worst case, but is usually much smaller in practice because most query
automata are sparse and length-deterministic. The intersection is not
materialized explicitly; following the workflow at the start of this
section, it is computed by a joint traversal of $D$ and $Q$ (lazy
depth-first search) or by a topological dynamic program with
bit-parallelism over the two automata. Critically, no query above visits
the dense $2^{n} \times 2^{n}$ matrix, even implicitly: the only objects
touched are the already-compressed automaton $D$ and the small query
automaton $Q$, so the asymptotic memory savings of Section~4 carry over
to query time. This makes the representation particularly attractive
for cryptanalytic workflows that revisit the DDT many times under
varying constraints.
